\chapter{وارون ماتریس}
\label{ch:inverse}

برای عدد $a\neq0$ وارون \enterm{Inverse} $\frac1a$ را داریم. برای ماتریس‌ها دو جور وارون پیدا می‌شود: وارون \emph{چپ} ($XA=I$) و وارون \emph{راست} ($AX=I$). در این فصل می‌بینیم هر کدام چه زمانی هست، چه چیزی درباره‌ی دستگاه $A\vx=\vb$ می‌گوید، و چرا برای ماتریس مربعی این دو یکی می‌شوند.

\section{وارون ماتریس مربعی}

\begin{definition}[وارون‌پذیر]
ماتریس $A_{n\times n}$ \term{وارون‌پذیر}{Invertible / Nonsingular} است اگر $B_{n\times n}$ وجود داشته باشد که $AB=BA=I_n$. آن‌گاه $B=A^{-1}$. ماتریس مربعی بدون وارون را \term{منفرد}{Singular} می‌نامیم.
\end{definition}

\begin{theorem}[یکتایی وارون]\label{thm:inv-unique}
وارون ماتریس مربعی یکتاست.
\end{theorem}
\begin{proof}
اگر $B$ و $C$ هر دو وارون $A$ باشند، $B=B(AC)=(BA)C=C$.
\end{proof}

\begin{theorem}\label{thm:square-inv-indep}
ماتریس مربعی وارون‌پذیر است اگر و تنها اگر ستون‌هایش مستقل باشند.
\end{theorem}
\begin{proof}
اگر $A$ وارون‌پذیر باشد و $A\vx=\zero$، آن‌گاه $\vx=A^{-1}A\vx=\zero$؛ پس ستون‌ها مستقل‌اند. برعکس، اگر ستون‌ها مستقل باشند، طبق قضیه‌ی \ref{thm:imt} (بندهای (ه) و (آ)) $A$ وارون‌پذیر است.
\end{proof}

\section{وارون چپ}

\begin{definition}[وارون چپ]
$X$ را یک \term{وارون چپ}{Left Inverse} $A_{m\times n}$ می‌نامیم اگر $XA=I_n$؛ پس $X$ ابعاد $n\times m$ دارد. اگر وجود داشته باشد، $A$ \term{چپ‌وارون‌پذیر}{Left-Invertible} است.
\end{definition}

وارون چپ لزوماً یکتا نیست. برای 
\[
A=\begin{bmatrix}-3&-4\\4&6\\1&1\end{bmatrix}
\]
هر دو ماتریس
\[
B=\frac19\begin{bmatrix}-11&-10&16\\7&8&-11\end{bmatrix},\qquad
C=\frac12\begin{bmatrix}0&-1&6\\0&1&-4\end{bmatrix}
\]
وارون چپ $A$ هستند (با ضرب مستقیم $BA=CA=I_2$ بررسی می‌شود).

\begin{theorem}[وارون چپ و استقلال ستون‌ها]\label{thm:left-inv}
ماتریس چپ‌وارون‌پذیر است اگر و تنها اگر ستون‌هایش مستقل باشند.
\end{theorem}
\begin{proof}
اگر $XA=I$ و $A\vx=\zero$، آن‌گاه $\vx=XA\vx=\zero$؛ پس ستون‌ها مستقل‌اند.

برعکس، اگر ستون‌های $A_{m\times n}$ مستقل باشند، طبق لم \ref{lem:ATA-null} داریم $\Nul(A\T A)=\Nul(A)=\{\zero\}$. ماتریس $A\T A$ مربعی ($n\times n$) با فضای پوچ \enterm{Null Space} صفر است؛ پس طبق قضیه‌ی \ref{thm:imt} وارون‌پذیر است و
$X=(A\T A)^{-1}A\T$
یک وارون چپ است: $XA=(A\T A)^{-1}(A\T A)=I$.
\end{proof}

\begin{note}
ماتریس \emph{بلند} ($m>n$) یا مربعی می‌تواند چپ‌وارون‌پذیر باشد؛ ماتریس \emph{پهن} ($m<n$) هرگز، چون بیش از $m$ ستون در $\R^m$ وابسته‌اند (قضیه‌ی \ref{thm:more-than-n}). به‌ویژه بردار ستونی ناصفر همیشه وارون چپ دارد (مثلاً $\frac{\va\T}{\norm{\va}^2}$) و بردار سطری هرگز ندارد.
\end{note}

\begin{example}
وارون چپ $\va=(1,0,3)\T$ را به دو روش بیابید.
\end{example}
\begin{solution}
(۱) $\frac1{a_i}\ve_i\T$ با $a_i\neq0$، مثلاً 
\[
\begin{bmatrix}1&0&0\end{bmatrix}.
\]
 \quad (۲) چون $\va\T\va=\norm{\va}^2$، ماتریس 
\[
\frac{\va\T}{\norm{\va}^2}=\frac1{10}\begin{bmatrix}1&0&3\end{bmatrix}
\]
هم وارون چپ است. اگر ستون‌های $A$ یکامتعامد \enterm{Orthonormal} باشند، $A\T A=I$ و $A\T$ وارون چپ $A$ است.
\end{solution}

\section{وارون راست}

\begin{definition}[وارون راست]
$X$ یک \term{وارون راست}{Right Inverse} $A$ است اگر $AX=I$. از $AX=I$ داریم $X\T A\T=I$؛ پس \textbf{$A$ راست‌وارون‌پذیر است اگر و تنها اگر $A\T$ چپ‌وارون‌پذیر باشد.}
\end{definition}

\begin{theorem}[وارون راست و استقلال سطرها]\label{thm:right-inv}
ماتریس راست‌وارون‌پذیر است اگر و تنها اگر سطرهایش مستقل باشند. چنین ماتریسی پهن یا مربعی است.
\end{theorem}
\begin{proof}
سطرهای $A$ ستون‌های $A\T$ هستند. طبق قضیه‌ی \ref{thm:left-inv}، $A\T$ چپ‌وارون‌پذیر است اگر و تنها اگر ستون‌هایش (سطرهای $A$) مستقل باشند.
\end{proof}

\section{وارون و حل دستگاه}

\begin{itemize}
  \item \textbf{وارون چپ:} اگر $XA=I$ و $A\vx=\vb$ جواب داشته باشد، $\vx=XA\vx=X\vb$. پس \emph{حداکثر یک} جواب هست، و تنها نامزد آن $X\vb$ است؛ اما باید با جایگذاری بررسی شود که واقعاً جواب است.
  \item \textbf{وارون راست:} اگر $AB=I$، آن‌گاه $A(B\vb)=\vb$؛ پس $\vx=B\vb$ برای \emph{هر} $\vb$ یک جواب است، بدون نیاز به بررسی سازگاری \enterm{Consistent}. اما یکتا بودنش معلوم نیست.
  \item \textbf{وارون دوطرفه:} $A\vx=\vb\iff\vx=A^{-1}\vb$؛ جواب همیشه هست و یکتاست.
\end{itemize}

\begin{example}
 
\[
A=\begin{bmatrix}-1&1\\0&-1\\2&1\end{bmatrix}
\]
 وارون چپ 
\[
X=\begin{bmatrix}1&2&1\\4&5&2\end{bmatrix}
\]
دارد ($XA=I$). دستگاه $A\vx=(-4,3,0)$ را بررسی کنید.
\end{example}
\begin{solution}
تنها نامزد $X\vb=(2,-1)$ است، اما $A(2,-1)=(-3,1,3)\neq\vb$. دستگاه ناسازگار \enterm{Inconsistent} است (کاهش سطری ستون آخر را محوری می‌کند). پس وارون چپ بدون بررسی جایگذاری کافی نیست. در مقابل، برای $A\T$ که وارون راست $X\T$ دارد، دستگاه $A\T\vx=(1,2)$ برای هر سمت راستی جواب دارد؛ مثلاً $B\T(1,2)=(\frac13,\frac23,-\frac23)$ و $C\T(1,2)=(0,\frac12,-1)$ دو جواب متفاوت‌اند.
\end{solution}

\section{وارون چپ و راست یکی می‌شوند}

\begin{theorem}\label{thm:left-right-square}
برای ماتریس مربعی $A$: اگر $A$ وارون چپ یا راست داشته باشد، وارون‌پذیر است و آن وارون همان $A^{-1}$ است.
\end{theorem}
\begin{proof}
اگر $XA=I$، طبق قضیه‌ی \ref{thm:left-inv} ستون‌های $A$ مستقل‌اند و طبق قضیه‌ی \ref{thm:square-inv-indep} $A$ وارون‌پذیر است؛ پس $X=XAA^{-1}=A^{-1}$. اگر $AX=I$، سطرها مستقل‌اند و $A\T$ وارون‌پذیر است؛ پس $A$ هم وارون‌پذیر است (با ترانهاده \enterm{Transpose} $(A\T)^{-1}$ وارون $A$ می‌شود) و $X=A^{-1}AX=A^{-1}$.
\end{proof}

\begin{center}
\begin{tikzpicture}[>=Latex, node distance=3cm, every node/.style={font=\small}]
  \node[draw, rounded corners, inner sep=5pt] (a) at (0,1.6) {\fa{چپ‌وارون‌پذیر}};
  \node[draw, rounded corners, inner sep=5pt] (b) at (7.2,1.6) {\fa{ستون‌های مستقل}};
  \node[draw, rounded corners, inner sep=5pt] (c) at (7.2,-1.2) {\fa{راست‌وارون‌پذیر}};
  \node[draw, rounded corners, inner sep=5pt] (d) at (0,-1.2) {\fa{سطرهای مستقل}};
  \draw[->, very thick, primary] (a)--(b) node[midway, above] {\scriptsize \fa{همیشه}};
  \draw[->, very thick, primary] (b)--(c) node[midway, right] {\scriptsize \fa{مربعی}};
  \draw[<-, very thick, primary] (d)--(c) node[midway, below] {\scriptsize \fa{همیشه}};
  \draw[->, very thick, primary] (d)--(a) node[midway, left] {\scriptsize \fa{مربعی}};
\end{tikzpicture}
\end{center}

\section{محاسبه‌ی وارون}

برای $n\times n$: ماتریس $[A\,|\,I_n]$ را با عملیات سطری \enterm{Row Operations} به RREF می‌بریم. اگر نتیجه $[I_n\,|\,B]$ شد، $B=A^{-1}$؛ اگر بخش چپ $I_n$ نشد، $A$ وارون ندارد. دلیل: $AX=I$ یعنی $n$ دستگاه $A\vx_j=\ve_j$ با ماتریس ضرایب \enterm{Coefficient Matrix} یکسان؛ پس یک کاهش سطری همه‌شان را حل می‌کند (روش فصل ۷ برای محاسبه‌ی وارون). یعنی $n\times n$ وارون‌پذیر است اگر و تنها اگر $\RREF(A)=I_n$.

\begin{example}
وارون $A=\begin{bmatrix}1&4\\-1&-3\end{bmatrix}$ را بیابید.
\end{example}
\begin{solution}
\[
\left[\begin{array}{cc|cc}1&4&1&0\\-1&-3&0&1\end{array}\right]\sim
\left[\begin{array}{cc|cc}1&0&-3&-4\\0&1&1&1\end{array}\right],
\qquad A^{-1}=\begin{bmatrix}-3&-4\\1&1\end{bmatrix}.
\]
\end{solution}

\begin{theorem}[ماتریس مقدماتی]\label{thm:elem-inv}
هر ماتریس مقدماتی \enterm{Elementary Matrix} وارون‌پذیر است و وارونش ماتریس مقدماتیِ همان نوع است که آن را به $I$ برمی‌گرداند (نتیجه‌ی \ref{cor:elem-invertible}).
\end{theorem}
مثلاً برای 
\[
E=\begin{bmatrix}1&0&0\\0&1&0\\-4&0&1\end{bmatrix}
\]
 (عمل $R_3\leftarrow R_3-4R_1$) وارون 
\[
E^{-1}=\begin{bmatrix}1&0&0\\0&1&0\\4&0&1\end{bmatrix}
\]
است.

\section{ویژگی‌های وارون}

\begin{theorem}\label{thm:inv-props}
اگر $A$ و $B$ وارون‌پذیر $n\times n$ باشند، $k$ عدد صحیح مثبت و $c\neq0$:
\begin{enumerate}
  \item $A^{-1}$ وارون‌پذیر است و $(A^{-1})^{-1}=A$.
  \item $(AB)^{-1}=B^{-1}A^{-1}$؛ و در حالت کلی $(A_1\cdots A_k)^{-1}=A_k^{-1}\cdots A_1^{-1}$.
  \item $(A^k)^{-1}=(A^{-1})^k=:A^{-k}$ و $A^kA^l=A^{k+l}$ برای هر $k,l\in\Z$ (با $A^0=I$).
  \item $(cA)^{-1}=\frac1cA^{-1}$.
  \item $A\T$ وارون‌پذیر است و $(A\T)^{-1}=(A^{-1})\T$.
\end{enumerate}
\end{theorem}
\begin{proof}
(۱) از $AA^{-1}=A^{-1}A=I$. (۲) $(AB)(B^{-1}A^{-1})=A(BB^{-1})A^{-1}=I$ و برعکس‌اش مشابه؛ بقیه با استقرا. (۳) از (۲) با $A_i=A$. (۴) $(cA)(\frac1cA^{-1})=I$. (۵) با ترانهاده گرفتن از $AA^{-1}=I$ و $A^{-1}A=I$: $(A^{-1})\T A\T=A\T(A^{-1})\T=I$.
\end{proof}

\begin{theorem}[دستگاه مربعی]
$A\vx=\vb$ ($n$ معادله، $n$ مجهول) \textbf{دقیقاً یک جواب} $C\vb$ دارد اگر و تنها اگر $A$ وارون‌پذیر و $C=A^{-1}$ باشد.
\end{theorem}
\begin{proof}
اگر $A$ وارون‌پذیر باشد، $\vx=A^{-1}\vb$ تنها جواب است. برعکس، اگر برای هر $\vb$ جواب یکتا باشد، $\Nul(A)=\{\zero\}$ (چون $\zero$ جواب $\vb=\zero$ است) و طبق قضیه‌ی \ref{thm:imt}، $A$ وارون‌پذیر است.
\end{proof}

\begin{theorem}[ماتریس مثلثی]\label{thm:triangular-inv}
ماتریس بالامثلثی یا پایین‌مثلثی با درایه‌های قطری ناصفر وارون‌پذیر است.
\end{theorem}
\begin{proof}
فرض کنید $L$ پایین‌مثلثی و $L\vx=\zero$. معادله‌ی اول $l_{11}x_1=0$ است و چون $l_{11}\neq0$، $x_1=0$. معادله‌ی دوم $l_{21}x_1+l_{22}x_2=0$ و در نتیجه $x_2=0$ و همین‌طور تا آخر (جایگذاری پیشرو). پس $\Nul(L)=\{\zero\}$ و طبق قضیه‌ی \ref{thm:imt} $L$ وارون‌پذیر است. برای بالامثلثی $R$، ماتریس $R\T$ پایین‌مثلثی با همان قطر است؛ پس وارون‌پذیر است و طبق بند ۵ قضیه‌ی \ref{thm:inv-props}، $R$ هم وارون‌پذیر است.
\end{proof}

\begin{note}
ماتریس تغییر پایه \enterm{Basis} همیشه وارون‌پذیر است، چون ستون‌هایش یک پایه‌اند و در نتیجه مستقل (قضیه‌ی \ref{thm:square-inv-indep}).
\end{note}

\section{رتبه و وارون}

\begin{theorem}\label{thm:rank-inverse}
(الف) اگر $M$ وارون‌پذیر باشد، $\rank(M)=\rank(M^{-1})=n$. (ب) اگر $A_{m\times n}$ و $B_{n\times n}$ وارون‌پذیر باشد، $\rank(AB)=\rank(A)$.
\end{theorem}
\begin{proof}
(ب) طبق قضیه‌ی \ref{thm:rank-AB}، $\rank(AB)\le\rank(A)$. از $A=(AB)B^{-1}$ هم $\rank(A)\le\rank(AB)$. (الف) رتبه‌ی \enterm{Rank} وارون‌پذیر $n$ است (قضیه‌ی \ref{thm:imt}) و $M^{-1}$ هم وارون‌پذیر است.
\end{proof}

\section{جمع‌بندی و تمرین}

\begin{summary}
\begin{itemize}
  \item چپ‌وارون‌پذیر $\iff$ ستون‌ها مستقل (وارون: $(A\T A)^{-1}A\T$)؛ راست‌وارون‌پذیر $\iff$ سطرها مستقل. وارون یک‌طرفه یکتا نیست.
  \item وارون چپ: حداکثر یک جواب (باید بررسی شود)؛ وارون راست: همیشه جواب هست.
  \item برای ماتریس مربعی وارون چپ و راست یکی‌اند و ستون‌ها و سطرها هر دو مستقل‌اند.
  \item $[A\,|\,I]\to[I\,|\,A^{-1}]$؛ $(AB)^{-1}=B^{-1}A^{-1}$ و $(A\T)^{-1}=(A^{-1})\T$.
  \item مثلثی با قطر ناصفر وارون‌پذیر است؛ $\rank(AB)=\rank(A)$ وقتی $B$ وارون‌پذیر باشد.
\end{itemize}
\end{summary}

\subsection*{تمرین‌ها}
\begin{enumerate}
  \item آیا $\begin{bmatrix}1&1\\0&1\\3&0\end{bmatrix}$ چپ‌وارون‌پذیر است؟ یک وارون چپ بیابید.
  \item وارون $\begin{bmatrix}2&1\\5&3\end{bmatrix}$ را به دو روش (فرمول $2\times2$ و گاوس-جردن \enterm{Gauss--Jordan Elimination}) بیابید.
  \item ثابت کنید اگر $A$ و $B$ مربعی باشند و $AB$ وارون‌پذیر باشد، هر دو وارون‌پذیرند.
\end{enumerate}

\subsection*{پاسخ‌های کوتاه}
\begin{enumerate}
  \item بله (دو ستون مستقل‌اند). $X=(A\T A)^{-1}A\T$ با $A\T A=\begin{bmatrix}10&1\\1&2\end{bmatrix}$ یک وارون چپ است؛ وارون چپ ساده‌تر: $\begin{bmatrix}1&-1&0\\0&1&0\end{bmatrix}$.
  \item $\begin{bmatrix}3&-1\\-5&2\end{bmatrix}$.
  \item $A(B(AB)^{-1})=I$ و $((AB)^{-1}A)B=I$؛ پس $A$ وارون راست و $B$ وارون چپ دارند و طبق قضیه‌ی \ref{thm:left-right-square} وارون‌پذیرند.
\end{enumerate}
